\section{Metoda}
\label{sec:method}

\subsection{Środowisko}
\label{sec:environment}

Model nie widzi danych zadania: zmienia i odczytuje ukryty stan tylko przez narzędzia, a każde wywołanie trafia do
zapisu przebiegu (rys.~\ref{fig:environment}).

\begin{figure}[htbp]
  \centering
  \begin{tikzpicture}[
      node distance=6mm and 7mm,
      box/.style={draw, rounded corners, align=center, font=\small, minimum height=10mm, minimum width=24mm},
      arr/.style={-{Stealth}, thick},
      both/.style={{Stealth}-{Stealth}, thick}]
    \node[box] (gen) {Generator\\zadań};
    \node[box, right=of gen] (sandbox) {Piaskownica\\(ukryty stan)};
    \node[box, right=of sandbox] (loop) {Pętla\\tool callingu};
    \node[box, right=of loop] (model) {Model LLM\\(przez adapter)};
    \node[box, below=of loop] (trace) {Zapis\\przebiegu};
    \node[box, left=of trace] (analyzer) {Analizator};
    \node[box, left=of analyzer] (refs) {Algorytmy\\wzorcowe};
    \draw[arr] (gen) -- (sandbox);
    \draw[both] (sandbox) -- (loop);
    \draw[both] (loop) -- (model);
    \draw[arr] (loop) -- (trace);
    \draw[arr] (trace) -- (analyzer);
    \draw[arr] (refs) -- (analyzer);
  \end{tikzpicture}
  \caption{Przepływ danych: wykonanie zadania (góra) i analiza po fakcie (dół).}
  \label{fig:environment}
\end{figure}

\TODO{generator z ziarnem, walidacja argumentów, limity, wstrzykiwanie błędów, adaptery dostawców.}

\subsection{Zadania}
\label{sec:tasks}

\begin{table}[htbp]
  \centering
  \caption{Rodziny zadań, narzędzia i granice teoretyczne.}
  \label{tab:tasks}
  \small
  \begin{tabular}{@{}p{2.8cm}p{3.4cm}p{4.2cm}p{3.6cm}@{}}
    \toprule
    Zadanie & Narzędzia & Wzorce & Granica \\
    \midrule
    Sortowanie $N$ elementów & \tool{compare(i, j)}, \tool{swap(i, j)} & bąbelkowe, przez wstawianie, przez wybór,
      szybkie, scalanie & $\lceil \log_2 N! \rceil$ porównań \\
    Wieże Hanoi, $n$ krążków & \tool{move(from, to)} & rekurencyjne & $2^n - 1$ ruchów \\
    Wyszukiwanie & \tool{read(index)} & binarne, liniowe, interpolacyjne & $\lceil \log_2 (N + 1) \rceil$ odczytów
      (binarne) \\
    \bottomrule
  \end{tabular}
\end{table}

\subsection{Rozpoznawanie algorytmu}
\label{sec:recognition}

\TODO{cechy przebiegu, odtworzenie wzorców na tych samych danych, odległość ciągów operacji, klasyfikacja: znany
algorytm, hybryda, strategia nowa; walidacja na syntetycznych przebiegach.}

\begin{algorithm}[htbp]
  \caption{Klasyfikacja przebiegu sortowania (szkic).}
  \label{alg:classify}
  \begin{algorithmic}[1]
    \Require zapis przebiegu $T$, dane wejściowe $x$, warianty wzorców $\mathcal{W}$, próg $\tau$
    \If{stan końcowy nie jest posortowany}
      \State \Return błąd i jego rodzaj
    \EndIf
    \ForAll{$w \in \mathcal{W}$}
      \State $d_w \gets$ odległość ciągu operacji $T$ od przebiegu $w(x)$
    \EndFor
    \If{$\min_w d_w < \tau$}
      \State \Return znany algorytm $\arg\min_w d_w$
    \ElsIf{fragmenty $T$ pasują do różnych wzorców}
      \State \Return hybryda
    \Else
      \State \Return strategia nowa
    \EndIf
  \end{algorithmic}
\end{algorithm}

\subsection{Miary}
\label{sec:measures}

\TODO{skuteczność; porównania wobec $\lceil \log_2 N! \rceil$ i wzorca, porównania zbędne na grafie relacji
elementów; zamiany wobec liczby inwersji i $N - c$, gdzie $c$ to liczba cykli permutacji; błędy; zgodność
deklaracji; tokeny, czas i koszt.}
